\documentclass[11pt,reqno]{amsart}
\usepackage[T1]{fontenc}\usepackage[utf8]{inputenc}\usepackage{lmodern}
\usepackage[a4paper,margin=24mm]{geometry}
\usepackage{amsmath,amssymb,mathtools,microtype,array,xurl}
\usepackage[hidelinks]{hyperref}
\newtheorem{theorem}{Theorem}\newtheorem{lemma}[theorem]{Lemma}
\newtheorem{corollary}[theorem]{Corollary}
\title[Primitive sextic cycles and square-source escape]{Primitive Sextic Cycles and Square-Source Escape\\for the Erd\H{o}s--Straus Selectors}
\author{The Clankers}\date{18 September 2026}
\begin{document}
\begin{abstract}
A canonical external-nonresidue factor graph can have a closed failed triangle
whose three actual stabilizer quotients all have order six. We supply the complete
factor data at a certified prime and prove the saturation and target exclusions.
Thus constant minimal sextic defects, including their nonzero cubic jets, do not
force canonical cycle closure. A separately marked square source does force an
exit from every canonical all-$3\bmod4$ triangle, with no favourable-factor
hypothesis. This is not a counterexample or a proof of Erd\H{o}s--Straus.
\end{abstract}\maketitle

Let $p$ be prime in one of the six hard classes
$1,121,169,289,361,529\pmod{840}$. Write $(\cdot/p)$ for its quadratic character.
At a prime $q<p$ with $(q/p)=-1$, the canonical source is
\[
 \sigma_q=\begin{cases}1&q\equiv3\pmod4,\\3&q\equiv1\pmod4,\end{cases}
 \quad R_q=\sigma_qq,\qquad A_q=(p+R_q)/4.
\]
Every prime factor $r\mid A_q$ with $(r/p)=-1$ is an edge $q\to r$, carrying
weight $v_r(A_q)$. These factors exist because $(A_q/p)=-1$, are below $p$, and
are distinct from $q$. This graph and its reciprocity rules are inherited [1].

For $a=\prod_t t^{e_t}$, retain the complete centered box and its integer counts
\[
 \beta_t\in[-e_t,e_t],\quad u=\prod_t t^{e_t+\beta_t},\quad
 v(\beta)=\prod_t t^{\beta_t}\pmod R,\quad
 \mu(g)=\#v^{-1}(g),\quad W=\operatorname{supp}\mu.
\]
The canonical channels are $R\mid4u+1$ (E) and $R\mid u+a$ (M). Their targets
are $-p^{-1}$ and $-1$; inversion retains the extra exterior orientation $-p$.
With $d=\gcd(a,u)$, set $(h,r,s)=(d^2/u,u/d,a/d)$. The exact returns are
\[
 E:(a,hs\kappa,phr\kappa),\quad \kappa=(pr+s)/R;
 \qquad M:(a,phs\lambda,phr\lambda),\quad\lambda=(r+s)/R. \tag{1}
\]
Their ordered divisor inverses are $a^2/(Ry-pa)$ and $pa^2/(Ry-pa)$, respectively.
These are the original Type I/II coordinates [2].

\begin{theorem}\label{tri}
For the prime
\[
 p=7510085481569082811681\equiv121\pmod{840},
\]
the canonical graph contains the closed component $31\to223\to307\to31$.
At each vertex its actual stabilizer is $K_q=G_q^6$, with
$G_q=(\mathbb Z/q\mathbb Z)^\times$, and both channels are empty. Every
inactive factor box fills $K_q$; each successor is the unique simple nonresidue
prime factor. All effective and full indices equal six.
\end{theorem}
\begin{proof}
Put
\[
 L_0=526211146410389771,\quad L_1=1165340089472441,\quad L_2=23500235136011.
\]
The complete factorizations are
\[
 \frac{p+31}{4}=2^4L_0\cdot223,\qquad
 \frac{p+223}{4}=2^7\cdot41L_1\cdot307,\qquad
 \frac{p+307}{4}=17\cdot19\cdot79\cdot101L_2\cdot31.       \tag{2}
\]
All factors and $p$ are prime by the executable complete-order certificate below.
The kernel orders are $5,37,51$. Generators are $2,2,171$.
At 31, the factor $2^4$ fills the kernel. At 223, the logs of $2,41$ are $1,11$;
$[-7,7]+11[-1,1]$ contains $[-18,18]$. At 307, the logs of $17,19,79,101$ are
$-17,7,-1,-2$. The last two fill $[-3,3]$; the other two give centres
\[
 -24,-17,-10,-7,0,7,10,17,24,
\]
whose consecutive gaps are at most seven. Their radius-three intervals contain
$[-27,27]$, so they fill the 51-element kernel. The remaining $L_i$ lie in their
respective kernels (residues $8,136,216$). Thus every inactive box is saturated.

If $r$ denotes the successor, its quotient class generates $C_6$, while $4\in K_q$.
Consequently $pK_q=rK_q$ and
$W_q=K_q\{r^{-1},1,r\}$. In successor coordinates the occupied classes are
$5,0,1$, and the middle/exterior target classes are $3,2,4$. The occupied set
has trivial quotient period, proving the actual stabilizer assertion.
For each factor $t\mid A_q$, reciprocity gives $(t/p)=(t/q)$. All inactive
factors are squares modulo $q$, and only $r$ is a nonresidue; its exponent is one.
This proves closure using all outgoing factors, not a selected subgraph.
\end{proof}

The integer coarse counts are $T_q(1,1,0,0,0,1)$, with $T_q=27,135,243$.
The complete source has $81+405+729=1215$ vectors, all serialized with $u,\beta$
and fine residues. The prime proof has 38 complete-order nodes and 42 trial leaves
below 10,000. Each internal node verifies a complete factorization of $n-1$ and
\[
 a^{n-1}=1\pmod n,\quad\gcd(a^{(n-1)/r}-1,n)=1\quad(r\mid n-1).
\]
These force the full order $n-1$ modulo every prime divisor of $n$, hence prove
primality. At the root, $p-1=2^5\cdot3\cdot5\cdot15646011419935589191$ and $a=11$.
The standard-library checker verifies every child. This is a finite, exact
certificate, not a probable-prime or infinite-family assertion.

A two-cycle on primes $3\pmod4$ would require both reciprocal quadratic symbols
to be $-1$, impossible. Thus the displayed cycle has minimum possible length,
although no least-prime claim is made. It is not an ES counterexample:
\[
 \begin{gathered}
 R=3,\quad a=1877521370392270702921,\quad u=71,\\
 (h,r,s,\kappa)=(71,1,26443962963271418351,2512176481510784743344)
 \end{gathered}
\]
returns the positive original E state through (1).

\paragraph{The odd component really survives.}
The primary Eisenstein factors of the three primes are
$-5-6\omega,-17-6\omega,-17-18\omega$, with $\omega^2+\omega+1=0$.
Using the convention and cubic symbol in [3, \S2, p.4; \S4, p.8], the exponents
of the primary and conjugate symbols are
\[
 H=\begin{pmatrix}*&1&0\\1&*&0\\0&0&*\end{pmatrix},\qquad
 J=\begin{pmatrix}*&0&1\\0&*&1\\2&2&*\end{pmatrix}.
\]
They obey $H=H^{\mathsf T}$ and $J=-J^{\mathsf T}$. Their sum is the rational
norm-symbol table. There is no reciprocity contradiction after the conjugate
prime is retained. The neighbour-square kernel elements are $16,196,17$, with
sixth roots $3,38,20$. The actual section carries $r^6=1,8,144$ also remain;
at 307 the extension $C_{306}\to C_6$ does not split.

For $\varepsilon^2=0$, the exact algebra isomorphism
\[
 \mathbb F_2[C_6]\longrightarrow
 \mathbb F_2[\varepsilon]/\varepsilon^2\times
 \mathbb F_4[\varepsilon]/\varepsilon^2,
 \quad F\longmapsto(F(1+\varepsilon),F(\omega(1+\varepsilon)))
\]
sends the reduced packet $1+X+X^5$ to $(1,\varepsilon)$.
An inverse uses $e_0=1+X^2+X^4$, $e_1=X^2+X^4$, $E=1+X^3$, $Z=X^4$:
\[
 (a+b\varepsilon,c+d\varepsilon)
 \longmapsto(a+bE)e_0+(c(Z)+Ed(Z))e_1.
\]
The cubic body is a supported zero and its jet is nonzero; none of this creates
an absent target. Squaring the packet gives $(1,0)$ but uses two active
occurrences, whereas the original active prime occurs only once.

\begin{theorem}[Square-nine escape]\label{nine}
For any canonical triangle of external nonresidue primes all $3\pmod4$, and
every vertex $q$ in it, the original source $(p+9q)/4$ is below $p$ and has a
nonresidue prime factor. Every such factor lies outside that triangle.
\end{theorem}
\begin{proof}
For its largest vertex $Q$, an incoming canonical edge gives
$4cQ=p+q_{\rm prev}<p+Q$, hence $3Q<p$. Therefore the square source is an
integer below $p$, and its character modulo $p$ is $-1$.
Its own vertex cannot divide it. The canonical successor is excluded by
\[
 \gcd((p+q)/4,(p+9q)/4)=\gcd((p+q)/4,2).
\]
The predecessor $s$ has $(q/s)=-1$, so $(s/q)=1$ by reciprocity. But every
nonresidue factor $t\mid(p+9q)/4$ has $(t/q)=-1$, again by reciprocity and the
square 9. Thus the predecessor is excluded too. Factorization supplies at least
one such $t$, proving the assertion without a prime-supply hypothesis.
\end{proof}

\begin{theorem}[A cardinality constraint on enlarged traps]
Let $C$ be a set of external nonresidue primes all $3\pmod4$. Suppose
$(2L-1)^2\max C<3p$ and $\min C>2L-2$. If every nonresidue factor of every
$(p+(2j+1)^2q)/4$, $q\in C$, $0\le j<L$, remains in $C$, then $|C|\ge2L+1$.
\end{theorem}
\begin{proof}
Every source supplies a factor. Different ports at $q$ have gcd
\[
 \gcd\left((p+(2i+1)^2q)/4,(j-i)(i+j+1)\right).
\]
Its prime divisors are at most $2L-2$, so no internal factor is shared. Each
vertex has at least $L$ successors. There are no self-edges or opposite pairs,
by reciprocity for square multipliers. Hence
$L|C|\le|C|(|C|-1)/2$.
\end{proof}

In a primitive sink strongly connected component the inactive factor $B$ in $A_q=Br$ is at least
four: $B=1$ would force $\operatorname{ord}_q(2)\mid6$; $B=2$ or $3$ would force
order three for that integer because its three-point box is a subgroup. These
allow only $q=3,7,13$, outside the domain. Thus its largest vertex obeys
$15Q<p$. The preceding theorem with $L=3$ proves that a primitive closed
component of at most six vertices has an exit at multiplier 9 or 25.

For the displayed large prime, the square-nine successors are respectively
$11$, $113306597729903$, and $268217338627467243373$, all outside the triangle.
Only the first square source is occupied (four E and four M states); the other
two are empty. Escape must not be identified with immediate selector success.
The next universal problem is to force an occupied original source after or
outside such extensions. Neither canonical reciprocity closure nor ES is proved.

\begin{thebibliography}{9}
\bibitem{1} Bryan, C. (2026). \emph{External-nonresidue factor cycle};
\emph{Primitive sextic defect chains and the neighbor-square recurrence}, \S\S2--5.
\texttt{chasebryan/centl}, revision \texttt{42508684c4386b575d7b52e763565345a76bcb9e}.
\bibitem{2} Elsholtz, C., \& Tao, T. (2013). Counting the number of solutions to
the Erd\H{o}s--Straus equation on unit fractions. \emph{J. Aust. Math. Soc., 94},
50--105. arXiv:1107.1010, \S2.
\bibitem{3} Damg\aa rd, I. B., \& Frandsen, G. S. (2003).
\emph{Efficient algorithms for gcd and cubic residuosity in the ring of Eisenstein integers}.
BRICS RS-03-8, \S2, p.4; \S4, p.8. DOI: 10.7146/brics.v10i8.21779.
\end{thebibliography}
\end{document}
